\documentclass{stex}

\libusepackage{naproche}
\libinput{preamble}

\begin{document}
\begin{smodule}{divisibility.ftl}

\importmodule[libraries/arithmetics]{multiplication-and-ordering.ftl}

\symdef{Divides}{\,\mid\,}
\symdef{NotDivides}{\,\nmid\,}

\symdecl*{divisible}
\symdecl*{divide}
\symdecl*{factor}
\symdecl*{divisor}
\symdecl*{trivial divisor}
\symdecl*{nontrivial divisor}
\symdecl*{divisibility of summands}

\begin{definition}[forthel,for={divisible,divide}]
  Let $n, m$ be natural numbers.
  $\emph{n \Divides m}$ iff there exists a natural number $k$ such that $n \Mul k \Eq m$.
  Let $m$ is \emph{divisible by $n$} stand for $n \Divides m$.
  Let $n$ \emph{divides $m$} stand for $n \Divides m$.
  Let $\emph{n \NotDivides m}$ stand for $\Not n \Divides m$.
\end{definition}

\begin{lemma}[forthel]
  Let $n, m$ be natural numbers.
  $n$ divides $m$ iff there exists a natural number $k$ such that $k \Mul n \Eq m$.
\end{lemma}

\begin{definition}[forthel,for={factor,divisor}]
  Let $n$ be a natural number.
  A \emph{factor of $n$} is a natural number that divides $n$.
  Let a \emph{divisor of $n$} stand for a factor of $n$.
\end{definition}

\begin{definition}[forthel,for=trivial divisor]
  Let $n$ be a natural number.
  A \emph{trivial divisor of $n$} is a divisor $m$ of $n$ such that $m \Eq \One$ or $m \Eq n$.
\end{definition}

\begin{definition}[forthel,for=nontrivial divisor]
  Let $n$ be a natural number.
  A \emph{nontrivial divisor of $n$} is a divisor $m$ of $n$ such that $m \NotEq \One$ and $m \NotEq n$.
\end{definition}

\begin{definition}[forthel]
  Let $n$ be a natural number.
  $n$ is \emph{composite} iff $n \Greater \One$ and $n$ has a nontrivial divisor.
\end{definition}

\begin{proposition}[forthel]
  Let $n$ be a natural number.
  Then $n \Divides \Zero$.
\end{proposition}
\begin{proof}[forthel]
  We have $n \Mul \Zero \Eq \Zero$.
  Hence $n \Divides \Zero$.
\end{proof}

\begin{proposition}[forthel]
  Let $n$ be a natural number.
  If $\Zero \Divides n$ then $n \Eq \Zero$.
\end{proposition}
\begin{proof}[forthel]
  Assume $\Zero \Divides n$.
  Consider a natural number $k$ such that $\Zero \Mul k \Eq n$.
  Then $n \Eq \Zero$.
\end{proof}

\begin{proposition}[forthel]
  Let $n$ be a natural number.
  Then $\One \Divides n$.
\end{proposition}
\begin{proof}[forthel]
  We have $\One \Mul n \Eq n$.
  Hence $\One \Divides n$.
\end{proof}

\begin{proposition}[forthel]
  Let $n$ be a natural number.
  Then $n \Divides n$.
\end{proposition}
\begin{proof}[forthel]
  We have $n \Mul \One \Eq n$.
  Hence $n \Divides n$.
\end{proof}

\begin{proposition}[forthel]
  Let $n$ be a natural number.
  If $n \Divides \One$ then $n \Eq \One$.
\end{proposition}
\begin{proof}[forthel]
  Assume $n \Divides \One$.
  Take a natural number $k$ such that $n \Mul k \Eq \One$.
  Suppose $n \NotEq \One$.
  Then $n \Less \One$ or $n \Greater \One$.

  \begin{case}{$n \Less \One$.}
    Then $n \Eq \Zero$.
    Hence $\Zero
      \Eq \Zero \Mul k
      \Eq n \Mul k
      \Eq \One$.
    Contradiction.
  \end{case}

  \begin{case}{$n \Greater \One$.}
    We have $k \NotEq \Zero$.
    Indeed if $k \Eq \Zero$ then
    $\One
      \Eq n \Mul k
      \Eq n \Mul \Zero
      \Eq \Zero$.
    Hence $k \Geq \One$.
    Take a positive natural number $l$ such that $n \Eq \One \Plus l$.
    Then $\One
      \Less \One \Plus l
      \Eq n
      \Eq n \Mul \One
      \Leq n \Mul k$.
    Hence $\One \Less n$.
    Contradiction.
  \end{case}
\end{proof}

\begin{proposition}[forthel]
  Let $n, m, k$ be natural numbers.
  If $n \Divides m$ then $n \Divides m \Mul k$.
\end{proposition}
\begin{proof}[forthel]
  Assume $n \Divides m$.
  Take $l \Elem \Naturals$ such that $n \Mul l \Eq m$.
  Then $n \Mul (l \Mul k)
    \Eq (n \Mul l) \Mul k
    \Eq m \Mul k$.
  Hence $n \Divides m \Mul k$.
\end{proof}

\begin{corollary}[forthel]
  Let $n, m, k$ be natural numbers.
  If $n \Divides m$ then $n \Divides k \Mul m$.
\end{corollary}

\begin{proposition}[forthel]
  Let $n, m, k$ be natural numbers.
  If $n \Divides m \Divides k$ then $n \Divides k$.
\end{proposition}
\begin{proof}[forthel]
  Assume $n \Divides m$ and $m \Divides k$.
  Take natural numbers $l,l'$ such that $n \Mul l \Eq m$ and $m \Mul l' \Eq k$.
  Then $n \Mul (l \Mul l')
    \Eq (n \Mul l) \Mul l'
    \Eq m \Mul l'
    \Eq k$.
  Hence $n \Divides k$.
\end{proof}

\begin{proposition}[forthel]
  Let $n, m$ be natural numbers such that $n \NotEq \Zero$.
  If $n \Divides m$ and $m \Divides n$ then $n \Eq m$.
\end{proposition}
\begin{proof}[forthel]
  Assume $n \Divides m$ and $m \Divides n$.
  Take natural numbers $k,k'$ such that $n \Mul k \Eq m$ and $m \Mul k' \Eq n$.
  Then $n
    \Eq m \Mul k'
    \Eq (n \Mul k) \Mul k'
    \Eq n \Mul (k \Mul k')$.
  Hence $k \Mul k' \Eq \One$.
  Thus $k \Eq \One \Eq k'$.
  Therefore $n \Eq m$.
\end{proof}

\begin{proposition}[forthel]
  Let $n, m, k$ be natural numbers.
  If $n \Divides m$ then $k \Mul n \Divides k \Mul m$.
\end{proposition}
\begin{proof}[forthel]
  Assume $n \Divides m$.
  Take a natural number $l$ such that $n \Mul l \Eq m$.
  Then $(k \Mul n) \Mul l
    \Eq k \Mul (n \Mul l)
    \Eq k \Mul m$.
  Hence $k \Mul n \Divides k \Mul m$.
\end{proof}

\begin{proposition}[forthel]
  Let $n, m, k$ be natural numbers.
  Assume $k \NotEq \Zero$.
  If $k \Mul n \Divides k \Mul m$ then $n \Divides m$.
\end{proposition}
\begin{proof}[forthel]
  Assume $k \Mul n \Divides k \Mul m$.
  Take a natural number $l$ such that $(k \Mul n) \Mul l \Eq k \Mul m$.
  Then $k \Mul (n \Mul l) \Eq k \Mul m$.
  Hence $n \Mul l \Eq m$ (by \sn{left-cancellability of multiplication}).
  Thus $n \Divides m$.
\end{proof}

\begin{proposition}[forthel]
  Let $n, m, k$ be natural numbers.
  If $k \Divides n$ and $k \Divides m$ then $k \Divides (n' \Mul n) \Plus (m' \Mul m)$
  for all natural numbers $n', m'$.
\end{proposition}
\begin{proof}[forthel]
  Assume $k \Divides n$ and $k \Divides m$.
  Let $n', m'$ be natural numbers.
  Take natural numbers $l,l'$ such that $k \Mul l \Eq n$ and $k \Mul l' \Eq m$.
  Then
  \[  k \Mul ((n' \Mul l) \Plus (m' \Mul l'))                \]
  \[    \Eq (k \Mul (n' \Mul l)) \Plus (k \Mul (m' \Mul l'))  \]
  \[    \Eq ((k \Mul n') \Mul l) \Plus ((k \Mul m') \Mul l')  \]
  \[    \Eq (n' \Mul (k \Mul l)) \Plus (m' \Mul (k \Mul l'))  \]
  \[    \Eq (n' \Mul n) \Plus (m' \Mul m).                      \]
\end{proof}

\begin{corollary}[forthel]
  Let $n, m, k$ be natural numbers.
  If $k \Divides n$ and $k \Divides m$ then $k \Divides n \Plus m$.
\end{corollary}
\begin{proof}[forthel]
  Assume $k \Divides n$ and $k \Divides m$.
  Take $n' \Eq \One$ and $m' \Eq \One$.
  Then $k \Divides (n' \Mul n) \Plus (m' \Mul m)$.
  $(n' \Mul n) \Plus (m' \Mul m) \Eq n \Plus m$.
  Hence $k \Divides n \Plus m$.
\end{proof}

\begin{proposition}[forthel,name=divisibility of summands]
  Let $n, m, k$ be natural numbers.
  If $k \Divides n$ and $k \Divides n \Plus m$ then $k \Divides m$.
\end{proposition}
\begin{proof}[forthel]
  Assume $k \Divides n$ and $k \Divides n \Plus m$.

  \begin{case}{$k \Eq \Zero$.} \end{case}

  \begin{case}{$k \NotEq \Zero$.}
    Take a natural number $l$ such that $n \Eq k \Mul l$.
    Take a natural number $l'$ such that $n \Plus m \Eq k \Mul l'$.
    Then $(k \Mul l) \Plus m \Eq k \Mul l'$.
    We have $l' \Geq l$.
    Indeed if $l' \Less l$ then
    $n \Plus m
      \Eq k \Mul l'
      \Less k \Mul l
      \Eq n$ (by \sn{preservation of ordering under left-multiplication}).
    Hence we can take a natural number $l''$ such that $l' \Eq l \Plus l''$.
    Then $(k \Mul l) \Plus m
      \Eq k \Mul l'
      \Eq k \Mul (l \Plus l'')
      \Eq (k \Mul l) \Plus (k \Mul l'')$.
    Indeed $k \Mul (l \Plus l'') \Eq (k \Mul l) \Plus (k \Mul l'')$ (by \sn{left-distributivity of multiplication over addition}).
    Thus $m \Eq (k \Mul l'')$ (by \sn{left-cancellability of addition}).
    Indeed $k \Mul l$ and $k \Mul l''$ are natural numbers.
    Therefore $k \Divides m$.
  \end{case}
\end{proof}

\begin{proposition}[forthel]
  Let $n, m$ be natural numbers such that $n, m \NotEq \Zero$.
  If $m \Divides n$ then $m \Leq n$.
\end{proposition}
\begin{proof}[forthel]
  Assume $m \Divides n$.
  Take a natural number $k$ such that $m \Mul k \Eq n$.
  If $k \Eq \Zero$ then
  $n
    \Eq m \Mul k
    \Eq m \Mul \Zero
    \Eq \Zero$.
  Thus $k \Geq \One$.
  Assume $m \Greater n$.
  Then $n
    \Eq m \Mul k
    \Geq m \Mul \One
    \Eq m
    \Greater n$.
  Hence $n \Greater n$.
  Contradiction.
\end{proof}

\end{smodule}
\end{document}
